\documentclass[11pt]{cdio}
%% Setup citation system
\usepackage[backend=biber, date=short, maxcitenames=2, bibencoding=utf8, style=apa]{biblatex}
\DeclareLanguageMapping{american}{american-apa}% needed for style=apa
%\renewcommand*{\bibfont}{\referencefont}
\setlength{\bibitemsep}{\baselineskip}
\setlength{\bibhang}{0pt}

\addbibresource{example-references.bib}% file with reference library

\graphicspath{{graphics/}{Graphics/}{./}}%You can store your images in these locations

\title{PAPER TITLE OF \underline{NOT MORE THAN 12 WORDS} (LATEX AND OVERLEAF VERSION)} 

\footer{Proceedings of the $21^\text{st}$ International CDIO Conference, hosted by\\Monash University, Melbourne, Australia, June 3--6, 2025}

\setuphyperref{}%enable clickable links
\begin{document}
\makecdiotitle{
  \author{Forename Surname, Forename Surname, etc. [Authors of SAME affiliation.]\\(Arial 11 pt, bold, center, spacing before 0 pt, spacing after 11pt)}
  \affiliation{Department, Institution [Affiliation of above authors.  Use APA style title case]}
  % https://apastyle.apa.org/style-grammar-guidelines/capitalization/title-case
  \author{Forename Surname, Forename Surname, etc. [Authors of SAME affiliation.]}
  \affiliation{Department, Institution [Affiliation of above authors.  Use APA style title case]}
}

\section{Abstract}
This section of the paper should be a \textbf{single paragraph abstract}
outlining the aims, scope and conclusion of the paper.
While no word limit is imposed, authors should aim for an abstract length of about \textbf{300 words} for full papers or about 50 words for shorter contributions.

\textbf{Please note that this template includes Microsoft Word instructions that are automatically implemented by this template.}
The text and pictures from the Word version are included to validate formatting but is unlikely to be relevant to \LaTeX{} users.
The template text should be removed before submission.

\section{Keywords}
Insert 4--6 keywords, separated by commas.
The last keyword must be ``Standards: \#'' containing a numerical list of particularly relevant CDIO standards (e.g. ``Standards: 1--6, 8, 10'').

\begin{center}
\fbox{
  \begin{minipage}{109mm}
    Note:
    \begin{itemize}
    \item All of the above must fit on the first page.
    \item Insert a page break before continuing with the rest of the paper
    \item Please remove this box before submitting.
    \end{itemize}
    \vspace{\baselineskip}
  \end{minipage}
}
\end{center}

%%%%%%% DO NOT REMOVE %%%%%%%%%%%%%%%
\newpage{}% Title, authors, abstract, and keywords go on the first page
%%%%%%%%%%%%%%%%%%%%%%%%%%%%%%%%%%%%%


\section{About CDIO papers and the CDIO paper template}
This section describes the main parameters of the CDIO paper template.
This template is itself formatted as a CDIO paper.

\textbf{CDIO conference papers are at most 10 pages, including the title page.
Longer manuscripts can be considered when appropriate; however, the authors are encouraged to write concisely, and the length of the paper must be justified by the substance of the contribution.}

The CDIO paper template uses a one-column format with left and right justification.
The page size is A4, with top, bottom, left, and right margins = 2.54 cm.
Gutter: 0 cm and Gutter position: Left.
The footer size is 1.25 cm.

The body text (style ``CDIO Body text'') has the following settings: Arial, 11 pt, regular, justified, single column, spacing before paragraph 0 pt, spacing after paragraph 11 pt.

Do not add an extra line before or after a body text paragraph --- the template takes care of that.
An exception is when adding an empty line or two (or a page break) to avoid an ugly page break.

The footer is in Arial, 9 pt, italic, centered, single line spacing, spacing before paragraph 0 pt, spacing after paragraph 0 pt.

\section{Formatting your paper}
This template is based on the ``Styles'' feature in Microsoft Word, where the formatting is automatically done for you by simply choosing the correct style.
See Figure~\ref{fig:word-cdio-paper_styles}


\begin{figure}[h]
  \centering
  \includegraphics[width=\textwidth]{word-cdio-paper_styles}
  \caption{CDIO paper styles}\label{fig:word-cdio-paper_styles}
\end{figure}

For this reason, we recommend typing directly in this file.
If you still wish to type your text in another program and copy the text into this file, make sure to choose the paste option “Keep Text Only” when pasting (shortcut Ctrl + Shift + V).
See Figure~\ref{fig:word-keep_text}.

\begin{figure}[t]
  \centering
  \includegraphics[width=30mm]{word-keep_text}
  \caption{Paste with ``Keep Text Only'' option}\label{fig:word-keep_text}
\end{figure}

After pasting your plain text, reformat the text by choosing the appropriate style.

We \textbf{strongly recommend using only} the pre-defined CDIO styles and settings in this document to format your paper.
Note the following guidelines:
\begin{itemize}
\item Do not modify styles or add new styles. 
\item Do not edit page margins, letter spacing, page size, etc. Note that the paper size is A4 standard (and not US letter size). 
\item Do not add page numbering.
\item Do not include footnotes or headers.
\item Paragraphs should not have any indents.
\item Dinot use landscape layout for pages (if needed: rotate the figure or table rather than the page).
\end{itemize}
If you have any questions, feel free to contact the CDIO Office at info@cdio.org.

\section{Headings}
Three levels of headings are provided in the template.

Do not add extra lines before or after a heading --- the template takes care of that.

Do not number the headings.

\section{Heading 1}
The primary heading uses the ``CDIO Heading 1'' style with the parameters: Arial, 11 pt, bold, all caps, left alignment, spacing before paragraph 22 pt, spacing after paragraph 11 pt.

\subsection{Heading 2}
The secondary heading uses the ``CDIO Heading 2'' style with the parameters: Arial, 11 pt, bold and italic, left alignment, spacing before paragraph: 0 pt, spacing after paragraph: 11 pt.

\subsubsection{Heading 3}

The tertiary heading uses the ``CDIO Heading 3'' style with the parameters: Arial, 11 pt, italic, left alignment, spacing before paragraph 0 pt, spacing after paragraph 11 pt.

\section{SPECIAL FORMATTING}

If you use Word’s caption feature, numbered list or other automatically numbered elements, you MUST convert them to plain text for the final submission.
This is to avoid runaway numbers when the papers are combined into the proceedings.

\subsection{Equations}

Equations may be inserted however you wish, if the following is followed:
\begin{itemize}
\item The equation is centered on the page. 
\item Equations are numbered in Arabic numbers in parenthesis, and the numbering is placed at the far right to the equation. 
\item There is roughly 11 pt of space before and after the equation. 
\end{itemize}
See an example below, showing an invisible-border 3-column table housing the equation (inserted using Microsoft Word’s built-in equation tool) in the center column and the equation number in the right-most column.
\begin{equation}\label{eq:1}
  a^{2} + b^{2} = c^{2}
\end{equation}
Question references should be referred to in the text in the form Eqn~(\ref{eq:1}).

Equations should be typed (e.g., images are not acceptable) and all symbols should be explained within the text of your paper. You may include a separate section detailing all nomenclature.

\subsection{Tables}
There is a degree of flexibility concerning tables.
You may choose how you wish to format the rows and columns, but please be consistent throughout your paper.
However, you must center the table on your page and insert a table caption using Word’s built-in “Insert Caption” tool.
Note that in accordance with the APA style, the table caption uses sentence case capitalization and does not end with a period.

Table captions use the style CDIO Caption with the following settings: Arial, 11 pt, centered alignment, single line spacing, spacing before paragraph (table caption) 0 pt, spacing after paragraph (table caption) 11 pt.
Place the table caption above the table.

We recommend to position tables at the top of pages.
Please ensure that a legible text size is chosen.
Tables should not span across multiple pages.
Look at Table~\ref{tab:gender_dist} for an example.
\begin{table}[t]
\centering
\caption{Gender-distribution between different levels of education\label{tab:gender_dist}}
\begin{tabular}{|l|l|l|l|}\hline
  &Male &Female &Total\\\hline
  Elementary &40 &60 &100\\\hline
  Junior High &60&40&100\\\hline
  High School &70&30&100\\\hline
\end{tabular}
\end{table}

\subsection{Figures}
When using figures, center the figure on the page, use the layout option ``In Line with Text'', and insert a figure caption using Word’s the build-in ``Insert Caption'' tool below the figure.
Note that in accordance with the APA style, the figure caption uses sentence case capitalization and does not end with a period.

Figure captions use the style CDIO Caption with the following settings: Arial, 11 pt, centered alignment, single line spacing, centered alignment, single line spacing, spacing before paragraph (figure caption) 0 pt.
spacing after paragraph (figure caption) 11 pt.
Place the figure caption below the figure.

Figures should appear as close to their first citation in the text as possible with consideration for a smooth layout.
We recommend to position figures at the top of pages.
Please ensure that figures are legible.

If photographs or images are included, high-quality originals should be used.

\begin{figure}[t]
  \centering
  \includegraphics[width=0.7\textwidth]{template-fig1}
   \caption{Entire caption for figure centered and below the illustration\label{fig:1}}
\end{figure}

\section{ACKNOWLEDGEMENTS}

CDIO paper should include an acknowledgement section, placed after the conclusions section and before the reference lists.
Use the ``CDIO Heading 1'' style for the heading.

In the acknowledgements section, Funding acknowledgements and other possible conflicts of interests should be listed, if applicable.
Provide the funding agency written out in full, followed by a grant number or other relevant details.
If no specific funding has been granted to support this work, please use the following sentence: ``The author(s) received no financial support for this work.''

You may also use the section to acknowledge individuals that helped you with the work but without being authors, such as interviewees, the reviewers etc.

A page break between the \textbf{ACKNOWLEDGEMENTS} and the \textbf{REFERENCES} section is not necessary.
\section{THE REFERENCE SECTION}
Use the APA format for all references as shown in the example below. If in doubt, consult \url{www.apastyle.org}\parencite{apa2018style}.

Use the pre-defined ``CDIO References'' style for your reference list. This style has the following main parameters: Arial, 10 pt, left alignment, spacing before paragraph 0 pt, spacing after paragraph 3 pt.

When citing a reference in the text, use Author/Date in parentheses~\parencite{biggs2011teaching} or \parencite{malmqvist2012application}.
If it reads better, refer directly to \textcite{lantada2012cdio} in the text.
(APA Rules say that you need to remove parenthesis if the citation is already in parenthesis like this:~\nptextcite{lantada2012cdio}).

We suggest that you use either a reference tool like endnote~\parencite{clarivate2018endnote}  or the built-in reference system in Word.
You can find a tutorial on how to use the built-in reference system in Word at \url{https://support.office.com/en-us/article/create-a-bibliography-citations-and-references-17686589-4824-4940-9c69-342c289fa2a5}\parencite{microsoft2018create}.
References should be made so that other people can find them; include all relevant information when refer to a webpage and the date it was accessed.

Please ensure that references cited in the main text are listed in the references list, and vice versa.

%\section{REFERENCES}
\setlength\bibitemsep{2pt}% Word template is single spaced, but it looks like there is extra space between items in the official implementation --foley

\printbibliography{}

\begin{center}
\fbox{
  \begin{minipage}{102mm}
      If the following \textbf{``BIOGRAPHICAL INFORMATION''} section fits on this page, insert it here.  If it does not \textit{fully} fit on this page, insert a page-break after the reference section.
  \end{minipage}
}
\end{center}

\newpage{}
\section{Biographical Information}
Include a 1-paragraph biography of the author(s) where the author’s name is in bold and italic font. For the corresponding author you must include institution, department, address and email. The authors must grant a Creative Commons version 4.0 license to reproduce the work and include the corresponding marking. See an example below.

\textbf{Reidar Lyng} is Associate Professor of university pedagogics at The Dept.
of Education and Lifelong Learning at NTNU, presently chairing the Centre for Science and Engineering Education Development (SEED) at NTNU in Trondheim.
He holds an MSc degree in Chemical Engineering and a PhD degree in Physical Chemistry.
He has more than 30 years’ experience of education development from NTNU and several Swedish universities.
His research and development interests are wide ranging and include the systemic interplay between teachers, students, and learning spaces.
Reidar is regional co-leader for the EU within CDIO.

\textbf{Geir Egil Dahle Øien} holds an MscEE degree, and a PhD degree in Telecommunications.
He became professor at NTNU in 2001, served as HoD 2002-2004, and as Dean 2009-2019.
He was FTS project manager 2019-2021.
Since Jan.
2022 he is back as professor at NTNU’s Dept.
of Electronic Systems.
Geir has extensive experience from periodic supervision of Nordic education institutions’ systematic quality work, He serves on the board of NORDTEK and leads the Portfolio Board for Enabling Technologies in The Research Council of Norway.
In 2022 he is guest researcher at Aarhus University in Denmark.

\subsection{Corresponding author}
\CorrespondingAuthorBlock{
Reidar Lyng\\
NTNU Norwegian University of Science and Technology\\
NTNU Center for Science \& Engineering Education Development SEED\\
Realfagbygget E1-125, Høgskoleringen 5\\
7491 Trondheim, NORWAY}
{Reidar.lyng@ntnu.no}

\begin{center}
\fbox{
  \begin{minipage}{102mm}
    Note:
    \begin{itemize}
    \item Insert a page break if you are going to add an appendix
    \end{itemize}
  \end{minipage}
}
\end{center}

\end{document}

%%% Local Variables:
%%% mode: latex
%%% TeX-master: t
%%% TeX-engine: xetex
%%% End:
